\documentclass[11pt,a4paper]{article}
\usepackage[utf8]{inputenc}
\usepackage[T1]{fontenc}
\usepackage[portuguese]{babel}
\usepackage{xcolor}
\usepackage{hyperref}
\usepackage{geometry}
\usepackage{listings}
\usepackage{tcolorbox}
\usepackage{graphicx}

\geometry{margin=2.5cm}

\definecolor{primary}{RGB}{41, 128, 185}
\definecolor{secondary}{RGB}{44, 62, 80}
\definecolor{codebg}{RGB}{240, 240, 240}
\definecolor{textcol}{RGB}{51, 51, 51}

\hypersetup{
    colorlinks=true,
    linkcolor=primary,
    filecolor=magenta,      
    urlcolor=primary,
}

\lstset{
    backgroundcolor=\color{codebg},
    basicstyle=\ttfamily\small\color{secondary},
    breaklines=true,
    frame=single,
    rulecolor=\color{codebg},
    keywordstyle=\color{primary}\bfseries,
    commentstyle=\color{gray}\itshape,
    stringstyle=\color{teal},
    showstringspaces=false
}

\title{\color{secondary}\Huge\textbf{Estratégias de Atualização e Prometheus}}
\author{\color{primary}DevOps com Kubernetes -- 2026}
\date{}

\begin{document}
\maketitle

\section*{\color{primary}O que você aprenderá nesta página}
\begin{itemize}
    \item Comparar estratégias de atualização
    \item Implantar software com lançamentos canário (\textit{canary releases})
    \item Usar o Prometheus para consultas personalizadas
\end{itemize}

No capítulo anterior, realizamos atualizações automatizadas com um pipeline de implantação. A atualização simplesmente funcionou, mas não fazemos ideia de como ela realmente ocorreu, a não ser pelo fato de que um \textit{pod} foi alterado. Vamos agora analisar como podemos tornar o processo de atualização mais seguro, ajudando-nos a alcançar um número maior de "noves" (\textit{nines}) na disponibilidade.

Existem múltiplas estratégias de atualização/implantação/lançamento. Vamos nos concentrar em duas delas:
\begin{itemize}
    \item Atualização contínua (\textit{Rolling update})
    \item Lançamento canário (\textit{Canary release})
\end{itemize}

Ambas as estratégias de atualização são projetadas para garantir que a aplicação continue funcionando durante e após uma atualização. Em vez de atualizar todos os \textit{pods} ao mesmo tempo, a ideia é atualizá-los um de cada vez e confirmar se a aplicação está funcionando.

\section*{\color{primary}Atualização contínua (\textit{Rolling update})}

Por padrão, o Kubernetes inicia uma atualização contínua (\textit{rolling update}) quando alteramos a imagem. Isso significa que cada \textit{pod} é atualizado sequencialmente. A atualização contínua é um excelente padrão, pois permite que a aplicação permaneça disponível durante o processo de atualização. Se decidirmos enviar uma imagem que não funciona, a atualização será interrompida automaticamente.

Preparei uma aplicação e imagens com 5 versões.
\begin{itemize}
    \item A versão com a tag v1 sempre funciona
    \item v2 nunca funciona
    \item v3 funciona 90\% das vezes
    \item v4 morrerá após 60 segundos, e
    \item v5 sempre funciona
\end{itemize}

A seguinte configuração pode ser usada para colocar a aplicação v1 em funcionamento:

\begin{lstlisting}[language=bash]
apiVersion: apps/v1
kind: Deployment
metadata:
  name: flaky-update-dep
spec:
  replicas: 4
  selector:
    matchLabels:
      app: flaky-update
  template:
    metadata:
      labels:
        app: flaky-update
    spec:
      containers:
        - name: flaky-update
          image: mluukkai/dwk-app8:v1

---

apiVersion: v1
kind: Service
metadata:
  name: flaky-update-svc
spec:
  selector:
    app: flaky-update
  ports:
    - protocol: TCP
      port: 80
      targetPort: 3541
  type: LoadBalancer
\end{lstlisting}

Vamos implantar a aplicação:

\begin{lstlisting}[language=bash]
$ kubectl apply -f deployment.yaml
  deployment.apps/flaky-update-dep created

$ kubectl get po
  NAME                                READY   STATUS    RESTARTS   AGE
  flaky-update-dep-7b5fd9ffc7-27cxt   1/1     Running   0          87s
  flaky-update-dep-7b5fd9ffc7-mp8vd   1/1     Running   0          88s
  flaky-update-dep-7b5fd9ffc7-m4smm   1/1     Running   0          87s
  flaky-update-dep-7b5fd9ffc7-nzl98   1/1     Running   0          88s
\end{lstlisting}

Como podemos ver, a aplicação está funcionando. Agora mude a tag para v2 e aplique.

\begin{lstlisting}[language=bash]
$ kubectl apply -f deployment.yaml
$ kubectl get po --watch 
  NAME                                READY   STATUS    RESTARTS   AGE
  flaky-update-dep-84f4c94bc7-dkqcw   1/1     Running   0          17s
  flaky-update-dep-84f4c94bc7-h7d2t   1/1     Running   0          9s
\end{lstlisting}

Você pode ver que a atualização contínua foi concluída, mas infelizmente, a aplicação não funciona mais. A aplicação está em execução, apenas existe um bug que a impede de funcionar corretamente. Como comunicamos esse mau funcionamento fora da aplicação? É aqui que entram as \textit{ReadinessProbes}.

Com uma \textit{ReadinessProbe}, o Kubernetes pode verificar se um \textit{pod} está pronto para processar requisições. A aplicação possui um \textit{endpoint} \texttt{/healthz} na porta 3541, e podemos usá-lo para testar a integridade. Ele simplesmente responderá com o código de status 500 se não estiver funcionando e 200 se estiver.

Vamos reverter a versão para a v1 também, para que possamos testar a atualização para a v2 novamente.

\begin{lstlisting}[language=bash]
apiVersion: apps/v1
kind: Deployment
metadata:
  name: flaky-update-dep
spec:
  replicas: 4
  selector:
    matchLabels:
      app: flaky-update
  template:
    metadata:
      labels:
        app: flaky-update
    spec:
      containers:
        - name: flaky-update
          image: mluukkai/dwk-app8:v1
          readinessProbe:
            initialDelaySeconds: 10 # Atraso inicial ate que a prontidao seja testada
            periodSeconds: 5 # Com que frequencia testar
            httpGet:
               path: /healthz
               port: 3541
\end{lstlisting}

Aqui, o \texttt{initialDelaySeconds} e o \texttt{periodSeconds} significam que a sonda é enviada 10 segundos depois que o contêiner estiver ativo e a cada 5 segundos depois disso. Agora, se alterarmos a tag para v2 e aplicarmos, três dos \textit{pods} estarão completamente funcionais. Um dos \textit{pods} v1 foi removido para dar lugar aos v2, mas como eles não funcionam, nunca ficam com status \textit{READY}, e a atualização não pode continuar.

\subsection*{\color{secondary}Continuando o curso com k3d}
Você pode continuar usando o Google Kubernetes Engine ou voltar a usar o k3d. Como nossas aplicações agora usam a Gateway API, você deve habilitá-la no cluster k3d.

\begin{tcolorbox}[colback=blue!5!white,colframe=blue!75!black,title=Exercício 4.1: Readiness probe]
Crie uma \textit{ReadinessProbe} para a aplicação Ping-pong. Ela deve estar pronta quando tiver uma conexão com o banco de dados. E crie outra \textit{ReadinessProbe} para a aplicação \textit{Log output}. Ela deve estar pronta quando puder receber dados da aplicação Ping-pong.
\end{tcolorbox}

\begin{tcolorbox}[colback=blue!5!white,colframe=blue!75!black,title=Exercício 4.2: O projeto, passo 21]
Crie as sondas (\textit{probes}) e o \textit{endpoint} necessários para O Projeto para garantir que ele esteja funcionando e conectado a um banco de dados. Adicione um botão à sua aplicação que pode ser usado para "quebrar" a aplicação.
\end{tcolorbox}

\begin{lstlisting}[language=bash]
let isHealthy = true
// ...
app.get("/healthz", (req, res) => {
  if (!isHealthy) {
    return res.status(500).json({ status: "unhealthy" })
  }
  return res.status(200).json({ status: "ok" })
})
\end{lstlisting}

\section*{\color{primary}Lançamento canário (\textit{Canary release})}

Com atualizações contínuas, ao incluir as Sondas (\textit{Probes}), poderíamos criar lançamentos sem tempo de inatividade para os usuários. Às vezes isso não é o suficiente e você precisa ser capaz de fazer um \textit{lançamento parcial} para alguns usuários e obter dados de uso e estatísticas para o próximo lançamento. \textit{Lançamento canário} é o termo usado para descrever uma estratégia de lançamento em que introduzimos um subconjunto dos usuários a uma nova versão da aplicação. Quando a confiança no novo lançamento aumenta, o número de usuários na nova versão pode ser aumentado até que a versão antiga não seja mais usada.

Usaremos o Argo Rollouts para testar um tipo de lançamento canário:

\begin{lstlisting}[language=bash]
$ kubectl create namespace argo-rollouts
$ kubectl apply -n argo-rollouts -f https://github.com/argoproj/argo-rollouts/releases/latest/download/install.yaml
\end{lstlisting}

O \textit{Rollout} substituirá nosso \textit{deployment} criado anteriormente e nos permitirá usar um novo campo:

\begin{lstlisting}[language=bash]
apiVersion: argoproj.io/v1alpha1
kind: Rollout
metadata:
  name: flaky-update-dep
spec:
  replicas: 4
  strategy:
    canary:
      steps:
      - setWeight: 25
      - pause:
          duration: 30s
      - setWeight: 50
      - pause:
          duration: 30s
# ... Sondas omitidas por brevidade
\end{lstlisting}

A estratégia acima moverá primeiro 25\% (\textit{setWeight}) dos \textit{pods} para uma nova versão, após o que esperará por 30 segundos, moverá para 50\% dos \textit{pods} e depois aguardará por 30 segundos até que cada \textit{pod} seja atualizado.

\begin{tcolorbox}[colback=blue!5!white,colframe=blue!75!black,title=Exercício 4.3: Prometheus]
Inicie agora o Prometheus com Helm e use \textit{port-forward} para acessar o site GUI. Escreva uma consulta que mostre o número de \textit{pods} criados por StatefulSets no \textit{namespace} \texttt{prometheus}.
\end{tcolorbox}

Com o novo \textit{Rollout} e \textit{AnalysisTemplate}, podemos testar com segurança a implantação de qualquer versão. O \textit{AnalysisTemplate} usará o Prometheus para consultar o estado do \textit{deployment}.

\begin{lstlisting}[language=bash]
apiVersion: argoproj.io/v1alpha1
kind: AnalysisTemplate
metadata:
  name: restart-rate
spec:
  metrics:
  - name: restart-rate
    initialDelay: 2m
    successCondition: result < 2
    provider:
      prometheus:
        address: http://kube-prometheus-stack-1602-prometheus.prometheus.svc.cluster.local:9090 # Nome DNS para o meu Prometheus, encontre o seu com kubectl describe svc ...
        query: |
          scalar(
            sum(kube_pod_container_status_restarts_total{namespace="default", container="flaky-update"}) -
            sum(kube_pod_container_status_restarts_total{namespace="default", container="flaky-update"} offset 2m)
          )
\end{lstlisting}

\begin{tcolorbox}[colback=blue!5!white,colframe=blue!75!black,title=Exercício 4.4: Seu canário]
Crie um \textit{AnalysisTemplate} para a aplicação Ping-pong que acompanhará o uso de CPU de todos os contêineres no \textit{namespace}.
\end{tcolorbox}

\section*{\color{primary}Outras estratégias de implantação}
O Kubernetes suporta a estratégia de Recriação (\textit{Recreate}), que derruba os \textit{pods} anteriores e substitui tudo pelo atualizado. O Argo Rollouts suporta a estratégia \textit{BlueGreen}, na qual uma nova versão é executada lado a lado com a nova, mas o tráfego é alternado entre as duas em um determinado ponto.

\begin{tcolorbox}[colback=blue!5!white,colframe=blue!75!black,title=Exercício 4.5: O projeto, passo 22]
Nossa aplicação de tarefas (\textit{todo}) poderia ter um campo "Concluído" (\textit{Done}) para tarefas que já foram feitas. Deve ser uma requisição PUT para \texttt{/todos/<id>}.
\end{tcolorbox}

\end{document}
